\documentclass[12pt,a4paper]{article}
\usepackage[utf8]{vietnam}
\usepackage{amsmath,amssymb}
\usepackage{tabularx}
\usepackage{tikz}
\usetikzlibrary{calc,angles,quotes,arrows.meta,patterns}
\usepackage{xcolor}
\usepackage{enumitem}
\usepackage{fancyhdr}
\usepackage{graphicx}
\usepackage[top=1.2cm, bottom=1.2cm, left=1.5cm, right=1.5cm]{geometry}
\usepackage[loigiai]{ex_test}
\SAOPTN{widthSA=2.6,heightSA=0.62}

\setlength{\headheight}{15pt}
\setlength{\parskip}{0pt}

\pagestyle{fancy}
\fancyhf{}
\rhead{\textit{Môn Toán 6 -- Hướng dẫn giải Đề số 15}}
\lhead{\textit{Bộ đề ôn tập Giữa học kì I}}
\cfoot{\thepage}

\begin{document}

\begin{center}
  \textbf{\Large HƯỚNG DẪN GIẢI CHI TIẾT VÀ BẢNG ĐÁP ÁN}\\
  \textbf{\large ĐỀ ÔN TẬP GIỮA HỌC KỲ I -- ĐỀ SỐ 15}\\
  \textit{Môn: Toán 6 -- Năm học 2026 -- 2027}
\end{center}
\vspace{2mm}

% --- BẢNG ĐÁP ÁN TRẮC NGHIỆM TÓM TẮT ---
\noindent\textbf{\large A. BẢNG ĐÁP ÁN TRẮC NGHIỆM TỔNG HỢP}
\vspace{1.5mm}

\noindent\textbf{1. Phần I (Mỗi câu trả lời đúng được 0,25 điểm):}
\begin{center}
\begin{tabular}{|c|c|c|c|c|c|c|c|c|c|c|c|c|}
\hline
\textbf{Câu} & \textbf{1} & \textbf{2} & \textbf{3} & \textbf{4} & \textbf{5} & \textbf{6} & \textbf{7} & \textbf{8} & \textbf{9} & \textbf{10} & \textbf{11} & \textbf{12} \\
\hline
\textbf{Đ/A} & \textbf{A} & \textbf{C} & \textbf{A} & \textbf{A} & \textbf{B} & \textbf{D} & \textbf{C} & \textbf{C} & \textbf{B} & \textbf{D} & \textbf{D} & \textbf{C} \\
\hline
\end{tabular}
\end{center}
\vspace{1mm}

\noindent\textbf{2. Phần II (Đúng 1 ý: 0,1đ; 2 ý: 0,25đ; 3 ý: 0,5đ; 4 ý: 1,0đ):}
\begin{center}
\begin{tabular}{|c|c|c|c|c||c|c|c|c|c|}
\hline
\textbf{Câu} & \textbf{a)} & \textbf{b)} & \textbf{c)} & \textbf{d)} & \textbf{Câu} & \textbf{a)} & \textbf{b)} & \textbf{c)} & \textbf{d)} \\
\hline
\textbf{Câu 1} & \textbf{Đ} & \textbf{S} & \textbf{S} & \textbf{Đ} & \textbf{Câu 2} & \textbf{Đ} & \textbf{Đ} & \textbf{Đ} & \textbf{S} \\
\hline
\end{tabular}
\end{center}
\vspace{1mm}

\noindent\textbf{3. Phần III (Mỗi câu trả lời đúng được 1,0 điểm):}
\begin{center}
\begin{tabular}{|c|c|c|}
\hline
\textbf{Câu} & \textbf{Câu 1} & \textbf{Câu 2} \\
\hline
\textbf{Đáp án} & \textbf{5} & \textbf{20} \\
\hline
\end{tabular}
\end{center}
\vspace{3mm}

% --- B. LỜI GIẢI CHI TIẾT ---
\noindent\textbf{\large B. LỜI GIẢI CHI TIẾT CÁC CÂU HỎI}
\vspace{2mm}

\noindent\textbf{\large PHẦN I. CÂU TRẮC NGHIỆM NHIỀU PHƯƠNG ÁN LỰA CHỌN}
\vspace{1mm}
\begin{ex}
Cho tập hợp $A = \{a; 5; b; 7\}$. Cách viết nào sau đây là đúng?
\choice
{\True $5 \in A$}
{$0 \in A$}
{$7 \notin A$}
{$a \notin A$}
\loigiai{
Số 5 là một phần tử thuộc tập hợp $A$ nên cách viết $5 \in A$ là đúng. Các cách viết còn lại sai vì: $0 \notin A$, $7 \in A$, $a \in A$.
}
\end{ex}

\begin{ex}
Cho hình vẽ sau:
\begin{center}
\includegraphics[width=0.45\linewidth]{crops/test_p57_xref3497.png}
\end{center}
Trong hình vẽ trên có:
\choice
{2 hình thang, 3 hình lục giác đều, 3 hình tam giác đều}
{2 hình thoi, 3 hình lục giác đều, 2 hình thang cân}
{\True 3 hình lục giác đều, 2 hình thang cân, 2 hình tam giác đều}
{2 hình tam giác đều, 1 hình lục giác đều, 4 hình thoi}
\loigiai{
Quan sát các hình cấu tạo nên hình vẽ:
- 3 hình lục giác đều (phần đầu, ngực, mông).
- 2 hình thang cân (phần thân giữa).
- 2 hình tam giác đều (phần chân).
}
\end{ex}

\begin{ex}
Các số La Mã $\text{XII; XXIII; XXIX}$ được đọc lần lượt là:
\choice
{\True 12; 23; 29}
{12; 23; 28}
{11; 23; 28}
{11; 23; 29}
\loigiai{
Ta có:
- $\text{XII} = 10 + 2 = 12$.
- $\text{XXIII} = 10 + 10 + 3 = 23$.
- $\text{XXIX} = 10 + 10 + (10 - 1) = 29$.
}
\end{ex}

\begin{ex}
Cho tập hợp $M = \{1; 2; 3; 4; 5; 6\}$. Khẳng định nào sau đây là đúng?
\choice
{\True $5 \in M$}
{$4 \notin M$}
{$8 \in M$}
{$6 \notin M$}
\loigiai{
Số 5 có mặt trong tập hợp $M$ nên khẳng định $5 \in M$ là đúng.
}
\end{ex}

\begin{ex}
Trong các khẳng định sau, khẳng định nào sai?
\choice
{$2^3 \cdot 2^4 = 2^7$}
{\True $201^0 = 0$}
{$5^7 : 5^5 = 25$}
{$2^4 = 16$}
\loigiai{
Quy ước mọi số tự nhiên khác 0 lũy thừa 0 đều bằng 1 ($a^0 = 1$ với $a \ne 0$). Do đó $201^0 = 1$, khẳng định "$201^0 = 0$" là sai.
}
\end{ex}

\begin{ex}
Giá trị của biểu thức $22 - 2 \cdot (3^2 + 1)$ bằng:
\choice
{8}
{140}
{200}
{\True 2}
\loigiai{
Thực hiện phép tính trong ngoặc trước:
$22 - 2 \cdot (3^2 + 1) = 22 - 2 \cdot (9 + 1) = 22 - 2 \cdot 10 = 22 - 20 = 2$.
}
\end{ex}

\begin{ex}
Số nào sau đây là số nguyên tố?
\choice
{6}
{1}
{\True 7}
{0}
\loigiai{
Số 7 là số tự nhiên lớn hơn 1 và chỉ có đúng 2 ước là 1 và 7, nên 7 là số nguyên tố.
}
\end{ex}

\begin{ex}
$\text{ƯCLN}(8, 20, 4)$ là:
\choice
{1}
{20}
{\True 4}
{8}
\loigiai{
Vì 8 và 20 đều chia hết cho 4 ($8 \ \vdots \ 4$ và $20 \ \vdots \ 4$), nên $\text{ƯCLN}(8, 20, 4) = 4$.
}
\end{ex}

\begin{ex}
Cho tập hợp $A$ gồm các số tự nhiên nhỏ hơn 5. Tập $A$ được viết là:
\choice
{$A = \{1; 2; 3; 4\}$}
{\True $A = \{0; 1; 2; 3; 4\}$}
{$A = \{1; 2; 3; 4; 5\}$}
{$A = \{0; 1; 2; 3; 4; 5\}$}
\loigiai{
Các số tự nhiên nhỏ hơn 5 là $0; 1; 2; 3; 4$. Vậy $A = \{0; 1; 2; 3; 4\}$.
}
\end{ex}

\begin{ex}
Cho hình vuông $ABCD$. Khẳng định nào sau đây là sai?
\choice
{Hình vuông $ABCD$ có bốn cạnh bằng nhau: $AB = BC = CD = DA$}
{Hình vuông $ABCD$ có bốn góc ở mỗi đỉnh $A, B, C, D$ bằng nhau}
{Hình vuông $ABCD$ có hai đường chéo bằng nhau: $AC = BD$}
{\True Hình vuông $ABCD$ có hai cặp cạnh đối song song: $AB$ và $BC$; $CD$ và $DA$}
\loigiai{
$AB$ và $BC$ là hai cạnh kề vuông góc với nhau, không phải hai cạnh đối song song. Hai cặp cạnh đối song song của hình vuông $ABCD$ phải là $AB \parallel CD$ và $AD \parallel BC$.
}
\end{ex}

\begin{ex}
Hình thang cân không có tính chất nào dưới đây?
\choice
{Hai cạnh bên bằng nhau}
{Hai đường chéo bằng nhau}
{Hai cạnh đáy song song}
{\True Hai góc kề một cạnh bên bằng nhau}
\loigiai{
Trong hình thang cân, hai góc kề một đáy bằng nhau, còn hai góc kề một cạnh bên có tổng bằng $180^\circ$ (bù nhau). Do đó khẳng định "Hai góc kề một cạnh bên bằng nhau" là sai.
}
\end{ex}

\begin{ex}
Trong các số sau: 2; 7; 11; 21, số nào là hợp số?
\choice
{2}
{7}
{\True 21}
{11}
\loigiai{
Số 21 chia hết cho 1; 3; 7; 21 (có nhiều hơn 2 ước) nên 21 là hợp số. Các số 2; 7; 11 đều là các số nguyên tố.
}
\end{ex}


\setcounter{ex}{0}
\noindent\textbf{\large PHẦN II. CÂU TRẮC NGHIỆM ĐÚNG SAI}
\vspace{1mm}
\begin{ex}
Cho tập hợp $X = \{x \in \mathbb{N} \mid 5 \le x < 10\}$. Xét tính đúng sai của các mệnh đề sau:
\choiceTF
{\True Tập hợp $X$ gồm 5 phần tử}
{Số 10 là một phần tử thuộc tập hợp $X$}
{Giá trị của biểu thức $3^2 + 4^2$ là một phần tử thuộc tập hợp $X$}
{\True Tổng tất cả các phần tử trong tập hợp $X$ bằng 35}
\loigiai{
\textbf{a)} Đúng. Các phần tử của $X$ là $5, 6, 7, 8, 9$, gồm đúng 5 phần tử.

\textbf{b)} Sai. Vì điều kiện $x < 10$ nên $10 \notin X$.

\textbf{c)} Sai. $3^2 + 4^2 = 9 + 16 = 25$. Vì $25 > 10$ nên $25 \notin X$.

\textbf{d)} Đúng. Tổng các phần tử là $5 + 6 + 7 + 8 + 9 = 35$.
}
\end{ex}

\begin{ex}
Cho hình tam giác đều $ABC$ có chu vi bằng $18\text{ cm}$. Xét tính đúng sai của các mệnh đề sau:
\choiceTF
{\True Độ dài mỗi cạnh của tam giác $ABC$ là $6\text{ cm}$}
{\True Ba góc tại các đỉnh $A, B, C$ của tam giác $ABC$ bằng nhau}
{\True Ghép 6 tam giác đều bằng nhau có cạnh $6\text{ cm}$ ta sẽ tạo thành một hình lục giác đều cạnh $6\text{ cm}$}
{Chu vi của hình lục giác đều được tạo thành ở trên là $48\text{ cm}$}
\loigiai{
\textbf{a)} Đúng. Tam giác đều $ABC$ có 3 cạnh bằng nhau nên mỗi cạnh có độ dài $18 : 3 = 6\text{ (cm)}$.

\textbf{b)} Đúng. Tam giác đều có 3 góc ở 3 đỉnh bằng nhau và cùng bằng $60^\circ$.

\textbf{c)} Đúng. Ghép 6 tam giác đều bằng nhau chung một đỉnh tạo thành hình lục giác đều có độ dài cạnh bằng cạnh tam giác đều ($6\text{ cm}$).

\textbf{d)} Sai. Hình lục giác đều có 6 cạnh bằng nhau nên chu vi là $6 \times 6 = 36\text{ (cm)}$, không phải $48\text{ cm}$.
}
\end{ex}


\setcounter{ex}{0}
\noindent\textbf{\large PHẦN III. CÂU HỎI TRẢ LỜI NGẮN}
\vspace{1mm}
\begin{ex}
Một trường Trung học cơ sở có 220 học sinh đi tham quan. Nhà trường cần thuê ít nhất bao nhiêu xe ô tô 45 chỗ ngồi để đủ chỗ cho tất cả học sinh? \shortans{5}
\loigiai{
Ta có: $220 : 45 = 4$ (dư 40). Nếu thuê 4 xe thì chỉ chở được 180 học sinh, còn dư 40 học sinh chưa có chỗ ngồi. Do đó nhà trường cần thuê ít nhất $4 + 1 = 5$ xe ô tô. Đáp số: 5.
}
\end{ex}

\begin{ex}
Hình chữ nhật có chiều dài và chiều rộng lần lượt là $6\text{ cm}$ và $4\text{ cm}$. Chu vi hình chữ nhật đó là bao nhiêu centimét? \shortans{20}
\loigiai{
Chu vi hình chữ nhật là: $(6 + 4) \times 2 = 20\text{ (cm)}$. Đáp số: 20.
}
\end{ex}


\setcounter{ex}{0}
\noindent\textbf{\large PHẦN IV. TỰ LUẬN (HƯỚNG DẪN GIẢI VÀ THANG ĐIỂM)}
\vspace{1mm}
\begin{ex} (2,0 điểm)
Thực hiện các phép tính sau (tính hợp lý nếu có thể):
\begin{enumerate}[label=\textbf{\alph*)}, leftmargin=0.8cm, itemsep=2pt]
\item (0,5 điểm) $270 : 3 : 5 - 4$
\item (0,5 điểm) $9^3 \cdot 2^3 - 4^2 : 2^4$
\item (0,5 điểm) $867 - (167 + 80)$
\item (0,5 điểm) $600 : \{450 : [450 - (4 \cdot 5^3 - 2^3 \cdot 5^2)]\}$
\end{enumerate}
\loigiai{
\textbf{a)} \textit (0,5 điểm): $270 : 3 : 5 - 4 = 90 : 5 - 4 = 18 - 4 = 14$.

\textbf{b)} \textit (0,5 điểm): $\begin{aligned}[t] 9^3 \cdot 2^3 - 4^2 : 2^4 &= (9 \cdot 2)^3 - (2^2)^2 : 2^4 \\ &= 18^3 - 2^4 : 2^4 \\ &= 5832 - 1 = 5831. \end{aligned}$

\textbf{c)} \textit (0,5 điểm): $867 - (167 + 80) = 867 - 167 - 80 = 700 - 80 = 620$.

\textbf{d)} \textit (0,5 điểm): $\begin{aligned}[t] 600 : \{450 : [450 - (4 \cdot 5^3 - 2^3 \cdot 5^2)]\} &= 600 : \{450 : [450 - (4 \cdot 125 - 8 \cdot 25)]\} \\ &= 600 : \{450 : [450 - (500 - 200)]\} \\ &= 600 : \{450 : [450 - 300]\} \\ &= 600 : \{450 : 150\} \\ &= 600 : 3 = 200. \end{aligned}$


}
\end{ex}
\vspace{1.5mm}

\begin{ex} (1,5 điểm)
Tìm số tự nhiên $x$, biết:
\begin{enumerate}[label=\textbf{\alph*)}, leftmargin=0.8cm, itemsep=2pt]
\item (0,25 điểm) $x - 15 = 35$
\item (0,25 điểm) $2(x + 15) = 5^2 \cdot 2$
\item (0,5 điểm) $128 - 3(x + 4) = 23$
\item (0,5 điểm) $x \ \vdots \ 15$ và $0 < x \le 60$
\end{enumerate}
\loigiai{
\textbf{a)} \textit (0,25 điểm): $x = 35 + 15 \Rightarrow x = 50$.

\textbf{b)} \textit (0,25 điểm): $2(x + 15) = 25 \cdot 2 = 50 \Rightarrow x + 15 = 50 : 2 = 25 \Rightarrow x = 25 - 15 = 10$.

\textbf{c)} \textit (0,5 điểm): $\begin{aligned}[t] 3(x + 4) &= 128 - 23 = 105 \\ x + 4 &= 105 : 3 = 35 \\ x &= 35 - 4 = 31. \end{aligned}$

\textbf{d)} \textit (0,5 điểm): Vì $x \ \vdots \ 15$ nên $x \in \text{B}(15) = \{0; 15; 30; 45; 60; 75; \dots\}$.
Kết hợp điều kiện $0 < x \le 60$, ta được $x \in \{15; 30; 45; 60\}$.


}
\end{ex}
\vspace{1.5mm}

\begin{ex} (1,0 điểm)
Một nền nhà hình chữ nhật của gia đình bác An có chiều dài là $12\text{ m}$, chiều rộng là $8\text{ m}$.
\begin{enumerate}[label=\textbf{\alph*)}, leftmargin=0.8cm, itemsep=2pt]
\item (0,5 điểm) Tính diện tích nền nhà hình chữ nhật đó.
\item (0,5 điểm) Bác An dự định lát nền bằng những viên gạch men hình vuông có cạnh là $40\text{ cm}$. Hỏi bác An phải dùng hết bao nhiêu viên gạch men để lát kín nền nhà đó? (biết diện tích mạch vữa là không đáng kể).
\end{enumerate}
\loigiai{
\textbf{a)} \textit (0,5 điểm): Diện tích nền nhà hình chữ nhật là:
$$S = 12 \times 8 = 96\text{ (m}^2\text{).}$$
Đáp số: $96\text{ m}^2$.

\textbf{b)} \textit (0,5 điểm): Đổi: $96\text{ m}^2 = 960\,000\text{ cm}^2$.
Diện tích một viên gạch men hình vuông là:
$$40 \times 40 = 1600\text{ (cm}^2\text{).}$$
Số viên gạch men bác An cần dùng để lát kín nền nhà là:
$$960\,000 : 1600 = 600\text{ (viên).}$$
Đáp số: 600 viên gạch.


}
\end{ex}
\vspace{1.5mm}

\begin{ex} (0,5 điểm)
Cho biểu thức $B = 2 + 21^{21} + 25^{120}$. Chứng minh rằng $B$ là hợp số.
\begin{enumerate}[label=\textbf{\alph*)}, leftmargin=0.8cm, itemsep=2pt]
\item (0,5 điểm) Chứng minh $B$ là hợp số.
\end{enumerate}
\loigiai{
\textbf{a)} \textit (0,5 điểm): Xét chữ số tận cùng của từng số hạng trong biểu thức $B$:
\begin{itemize}[leftmargin=0.6cm, itemsep=1pt]
\item Số 2 có chữ số tận cùng là 2.
\item $21^{21}$ có chữ số tận cùng là 1 (vì số có chữ số tận cùng là 1 khi nâng lên lũy thừa bất kỳ luôn tận cùng là 1).
\item $25^{120}$ có chữ số tận cùng là 5 (vì số có chữ số tận cùng là 5 khi nâng lên lũy thừa bất kỳ luôn tận cùng là 5).
\end{itemize}
Chữ số tận cùng của biểu thức $B$ là chữ số tận cùng của tổng $2 + 1 + 5 = 8$.
Vì $B$ có chữ số tận cùng là 8 nên $B \ \vdots \ 2$.
Mặt khác, $B > 2$ và $B \ \vdots \ 2$ nên ngoài 1 và chính nó, $B$ còn có ước là 2.
Vậy $B$ là hợp số.


}
\end{ex}
\vspace{1.5mm}


\begin{center}
\rule{4cm}{0.4pt}\\
\textit{--- HẾT ---}
\end{center}
\end{document}
